\chapter{Input and output}\label{ch-transput}

The programs so far used \code|print| and \code|read| on the standard
output and the standard input. This chapter describes input and output
in general: how a program reads and writes files, reads from and writes to
strings, reacts to the end of the input and to bad data, converts numbers
to text, and lays out its output precisely with formats. The Report
calls input and output together \rr{transput}.

The transput of the Report is a large part of its standard library, and
the implementations provide it to different degrees. Algol 68 Genie
provides nearly all of it; Section~\ref{sec-transput-genie} lists what it
lacks. For ga68 and a68j, the module \code|Genie_Transput| provides the
unformatted part used in this chapter; formatted and binary transput
are not available there, and the examples of those sections were run
with Algol 68 Genie.

\section{Books, channels and files}\label{sec-books}

The Report distinguishes three things that current languages usually do
not separate:
\begin{itemize}
\item A \rr{book} is the data itself, such as a file on disk, identified
  by a name. Its text is divided into pages, the pages into lines, and
  the lines into characters, and a position in a book is given by a page,
  a line and a character number.
\item A \rr{channel} describes a kind of device and what can be done with
  it: \code|stand_in_channel| for reading, \code|stand_out_channel| for
  writing, \code|stand_back_channel| for both.
\item A \rr{file} (here of type \code|File|) is what a program uses to
  read or write a book through a channel: a file handle or stream, in
  today's terms. It records the current position in the book, the
  procedures to call when something unusual happens, and a few settings.
\end{itemize}

The procedure \code|open (f, name, channel)| connects the file
\code|f| with an existing book; \code|establish (f, name, channel)|
creates a new one, and \code|create (f, channel)| creates a new book
without a name, a temporary file. These procedures return 0 on success.
\code|close (f)| ends the connection, \code|scratch (f)| ends it and
deletes the book, and \code|reset (f)| returns to the start of the
book. The files \code|stand_in| and \code|stand_out| are open from the
start; \code|print (x)| is short for \code|put (stand_out, x)|, and
\code|read (x)| for \code|get (stand_in, x)|.

\begin{implnote}
In ga68 and a68j the type is called \code|File| and the channels and
procedures are those of \code|Genie_Transput|, since the reserved word
\code|file| of ga68 does not yet denote a usable type. In the Report and
in Algol 68 Genie, the type is \code|file|. In the Report,
\code|establish| has three more parameters, the number of pages, lines
per page and characters per line of the new book; the implementations
leave them out.
\end{implnote}

\section{Reading and writing files}\label{sec-get-put}

The procedures \code|put| and \code|get| write and read a list of values,
like \code|print| and \code|read|, but on a given file. The following
program writes five lines to a new file and then reads them back:

\example{ch09/files}

The reading loop has no condition: it ends when the input is exhausted.
The procedure \code|on_logical_file_end| installs a procedure that the
transput calls when a \code|get| reaches the end of the book; here it
jumps out of the loop to the label \code|done|. Such procedures are the
Report's \rr{event routines}; today they would be called callbacks or
handlers. Section~\ref{sec-events} describes them in general.

The predefined label \code|stop| ends the program; \code|goto stop|, or
just \code|stop|, can be written anywhere.

\begin{implnote}
In Algol 68 Genie, \code|establish| returns 0 even if the file already
exists; the error comes with the first \code|put| (``file exists''). The
program above therefore fails when it is run a second time in the same
directory. \code|open| with \code|stand_out_channel| opens an existing
file for writing and empties it. A program that should work in both cases
tries \code|open| first and calls \code|establish| only if that fails.
\end{implnote}

\section{Reading the standard input}\label{sec-stand-in}

The same technique reads all of the standard input, however much there
is. Numbers are separated by blanks and line breaks. One detail matters
here: the event routine jumps to a label of the program's block, so it
lives only as long as that block (Section~\ref{sec-proc-lifetimes}),
whereas \code|stand_in| lives as long as the whole program. Installing
the routine in \code|stand_in| itself would break the lifetime rule; the
Report uses a case like this as its example. The program therefore
copies \code|stand_in| into a local file variable and reads from that:

\example{ch09/sum-input}

\noindent with the input
\lstinputlisting[style=output]{examples/ch09/sum-input.in}

The copy \code|input| reads the same data as \code|stand_in|: a
\code|File| value contains references to the book and to the current
position, which the copy shares, and its own event routines, which it
does not share. The copy can therefore be given an event routine that
lives only as long as the block, while reading continues where
\code|stand_in| stands.

A number is read up to the first character that cannot belong to it:
\code|12abc| yields 12 and leaves \code|abc| for the next read, and
\code|1,024| yields 1. Leading blanks and line breaks are skipped. A
character is read as it is, a blank included; at the end of a line, the
line break is skipped and the first character of the next line is read.

A \code|string| is read up to the end of the line; the line break itself
is not read, and the next \code|get| of a string would yield an empty
string at once. It is skipped with \code|newline|, which, like
\code|space|, \code|backspace| and \code|newpage|, can appear in the list
of \code|get| as well as of \code|put|. On input they skip the rest of
the line, one character, and so on; on output they start a new line,
move one position, and start a new page. A program that processes its
input line by line reads a line, processes it, and then skips the line
break:

\example{ch09/lines}

\noindent The input was:
\lstinputlisting[style=output]{examples/ch09/lines.in}

\begin{pitfall}
The shorter \code|get (input, (line, newline))| looks equivalent, but
it loses the last line if the input does not end with a line break:
\code|newline| then reaches the end of the input, the event routine
jumps away, and the line that was just read is never printed. Text
editors and programs do not always write a final line break.
\end{pitfall}

The layout procedures are ordinary procedures of type
\code|proc (ref File) void|. They can be called directly, as
\code|newline (stand_out)|, and a program can write its own procedure of
this type and put it in a list, where \code|put| and \code|get| call it
with the file:

\example{ch09/layout}

\section{Events}\label{sec-events}

An event routine is installed in a file for one kind of event. It
receives the file and returns a \code|bool|. If it returns \code|true|,
it reports that it has dealt with the event and the transput continues;
if it returns \code|false|, the default action follows. A routine can
also jump out, as in the examples above.

\begin{center}
\begin{tabular}{lll}
\toprule
Installed with & Called when & Default action \\
\midrule
\code|on_logical_file_end| & the data of the book are exhausted & error \\
\code|on_physical_file_end| & the medium is full or exhausted & error \\
\code|on_line_end| & a line is exhausted & \code|newline| \\
\code|on_page_end| & a page is exhausted & \code|newpage| \\
\code|on_format_end| & a format is exhausted (Section~\ref{sec-format-list}) & start it again \\
\code|on_value_error| & a value cannot be read or written & error \\
\code|on_char_error| & a character cannot be converted & error \\
\bottomrule
\end{tabular}
\end{center}

The Report leaves the program undefined where the table says ``error'';
the implementations stop it with a message. \code|Genie_Transput|
provides the logical end of the file and the value error. Algol 68 Genie
provides all except \code|on_char_error|; a character that cannot be
converted causes a value error.

A routine that returns \code|true| must really have dealt with the
event, because the transput then continues at the same place. With
Algol 68 Genie and \code|Genie_Transput|, a value error on input occurs
when the text is not a number. (The Report calls the char error routine
for a character that cannot belong to a number, and the value error
routine for a number that cannot be converted.) If the routine only
counts the error, the bad text is still there, and the next number fails
in the same way:

\example{ch09/value-error-norepair}

\noindent prints \code|12 -1 -1 -1, errors: 3|: one bad word spoiled all the
following values. The routine must remove the cause, here by reading the
bad word as a string:

\example{ch09/value-error}

\noindent which prints \code|12 -1 45 67, errors: 1|. In the same way, a
logical file end routine that returns \code|true| without having made
more data available is called again, according to the Report; the
implementations call it once and continue, and the value that was to be
read is then not meaningful. Such a routine should jump away, as in the examples above, or
return \code|false|.

\section{Strings as files}\label{sec-associate}

The procedure \code|associate (f, s)| connects a file with a string
variable instead of a book. Reading then takes the characters from the
string, and writing adds to it. This converts between numbers and
their textual form with all the facilities of the transput:

\example{ch09/associate}

\begin{implnote}
In the Report, \code|associate| connects a file with an array of arrays
of arrays of characters (pages, lines, characters) of fixed size, and
writing overwrites it from the start. Algol 68 Genie and
\code|Genie_Transput| use a \code|ref string| instead and append to it,
which is what programs need in practice.
\end{implnote}

\section{Terminators and other settings}\label{sec-make-term}

A string read with \code|get| ends at the end of the line. The
procedure \code|make_term (f, s)| makes it end also at every character of
the string \code|s|; the terminating character itself is not read. The
complete program at the end of this chapter uses it to read a name that
ends at a blank.

A file can be asked what it allows: \code|get_possible (f)|,
\code|put_possible (f)|, \code|bin_possible (f)|,
\code|reset_possible (f)| and \code|set_possible (f)| return a
\code|bool|. After a file has been written, \code|reset| returns to the
start so that it can be read; reading and writing cannot simply
alternate.

The Report provides further procedures, which Algol 68 Genie does not
have: the current position (\code|char_number|, \code|line_number|,
\code|page_number|) and the setting of it (\code|set_char_number| and a
\code|set| with page, line and character), a new name for a book
(\code|reidf|), and a translation table for characters
(\code|make_conv|). A program that needs the current line number counts
the lines itself. \code|Genie_Transput| provides, besides
\code|make_term|, the functions \code|eof| and \code|eoln| of Algol 68
Genie, which test for the end of the file and of the line.

\section{Printing arrays and records}\label{sec-straightening}

\code|put| and \code|print| accept arrays and records, not only single
values, provided that they contain no references, functions or unions. They are broken down into their elements
and fields,
which are printed one after the other; the Report calls this
\rr{straightening}. An array of records thus prints the fields of every
element:

\example{ch09/straightening}

\code|get| straightens in the same way: \code|get (f, m)| for a
\code|[2,3] int| variable \code|m| reads six numbers into it, row by row.
For input the Report is stricter: apart from a \code|string| variable
itself, the variable may contain no resizable arrays, so a record with a
\code|string| field or an array of strings cannot be read with one
\code|get|.

\begin{implnote}
\code|Genie_Transput| straightens arrays of the primitive types but not
records (Section~\ref{sec-print-data}), and its \code|get| does not
accept arrays; this example was run with Algol 68 Genie, whose
\code|int| has 64 bits and is therefore printed in a wider field. Algol
68 Genie also reads a record with \code|string| fields with one
\code|get|.
\end{implnote}

\section{Converting numbers to text}\label{sec-conversions}

The procedures \code|whole|, \code|fixed| and \code|float| of
Section~\ref{sec-io} convert a number into a string of a given width,
which is the usual way to lay out numbers without formats:
\begin{itemize}
\item \code|whole (i, w)| writes an integer in a field of width
  $|w|$. If \code|w| is positive, a sign is always written; if it is
  negative, only a minus sign. \code|w = 0| gives the shortest string.
\item \code|fixed (x, w, d)| writes a real with \code|d| digits after
  the point, in a field of width $|w|$, with the same rules for the sign.
\item \code|float (x, w, d, e)| writes a real in scientific notation with
  \code|d| digits after the point and an exponent of \code|e| digits.
\end{itemize}

\example{ch09/conversions}

\noindent prints
\outputof{ch09/conversions}

A number that does not fit is not given a wider field, unlike C's
\texttt{printf}: \code|whole (123456, 3)| yields three asterisks.
\code|fixed| first gives up digits after the point, so that
\code|fixed (3.14159, 6, 4)| yields \code|+3.142|; only when the integral
part does not fit either does it yield asterisks. \code|fixed| writes no
zero before the point of a number below 1 if the width is 0
(\code|.0150|). \code|float| chooses the exponent so that as many digits
as possible stand before the point.

\section{Formatted output}\label{sec-formats}

For complete control over the layout, the Report provides
\term{formats}, a small language of their own, comparable to the format
strings of \texttt{printf} in C or of Fortran's \texttt{FORMAT}
statement. A format text is written between dollar signs. The procedure
\code|printf| takes a list that begins with a format and continues with
the values; \code|putf| does the same for a given file, and
\code|readf| and \code|getf| read according to a format.

\example{ch09/formats}

A format consists of \term{pictures}, each of which describes how one
value is printed, separated by commas. A picture is made of
\term{frames}:

\begin{center}
\begin{tabular}{ll}
\toprule
Frame & Meaning \\
\midrule
\code|d| & a digit \\
\code|z| & a digit, printed as a blank if it is a leading zero \\
\code|+|, \code|-| & a sign: always, or only if negative \\
\code|.| & the decimal point \\
\code|e| & the exponent marker of a floating-point number \\
\code|a| & a character \\
\code|b| & a truth value (\code|T| or \code|F|) \\
\code|g| & a value in its standard layout, as \code|print| would print it \\
\bottomrule
\end{tabular}
\end{center}

A number before a frame repeats it: \code|5d| is five digits and
\code|3zd| is three suppressible digits followed by one that is always
printed. Between the frames, \term{insertions} add text: a string in
quotes is printed as it is, \code|x| prints a blank, \code|l| ends the
line and \code|p| the page, and \code|k| moves to a given column. They
too can be repeated, as in \code|2x| or \code|3"-"|.

\subsection*{One picture, one value}

Everything between two commas is one picture, and one picture prints
one value, however many frames it has:

\example{ch09/pictures}

\noindent prints
\outputof{ch09/pictures}

In the first format, \code|2d ", b = " 2d| is a single picture of four
digits with some text in the middle, and it prints the single value 1 as
\code|0001| with the text inserted. This is easy to get wrong and hard to
see; when a format prints something strange, the commas are the first
thing to check.

\subsection*{Where the sign goes}

A sign frame before the digits puts the sign at the left edge of the
field. A sign frame after the \code|z| frames lets it move to just before
the first printed digit:

\example{ch09/signs}

\noindent prints
\outputof{ch09/signs}

A negative value in a picture without a sign frame is a value error.

\subsection*{Choices, general patterns and other kinds of values}

A truth value can be printed with two strings of the program's choice,
\code|b("yes", "no")|, and an integer between 1 and $n$ can select one of
$n$ strings with \code|c("one", "two", ...)|. Both make it easy to print
plurals and names:

\example{ch09/choice}

\noindent prints
\outputof{ch09/choice}

The general pattern \code|g| prints a value as \code|print| would. With
parameters, it prints a number like the conversion procedures of
Section~\ref{sec-conversions}: \code|g(w)| like \code|whole|,
\code|g(w, d)| like \code|fixed|, and \code|g(w, d, e)| like
\code|float|. The parameters can be any integer expressions. A
\code|bits| value is printed in base 2, 4, 8 or 16 with a picture such as
\code|2r8d| or \code|16r4d|:

\example{ch09/general}

\noindent prints
\outputof{ch09/general}

The first line shows the standard layout of a \code|real| in Algol 68
Genie, whose \code|real| has 64 bits.

A string is printed with \code|a| frames, one per character, and the
number of frames must equal the length of the string; insertions between
the frames are allowed:

\example{ch09/strings-fmt}

\noindent prints \code|J. Smith| and \code|[Algol]|, and then stops with
an error, since \code|"Al"| has only two characters for five frames. A
string of variable length is printed with \code|g| or with the
computed repetition of the next section.

A part of a format can also be chosen when the program runs. The format pattern
\code|f(...)| contains an expression that yields a format, and inserts
that format in its place:

\example{ch09/format-f}

\noindent prints
\outputof{ch09/format-f}

\begin{pitfall}
Two string insertions with nothing between them are one string: as in a
string literal, a doubled quote stands for a quote character, so
\code|"Smith""Jones"| prints \texttt{Smith"Jones}, as the fourth line
shows. A repetition count, even \code|1|, keeps them apart. Algol 68
Genie also accepts a space between them, the Report does not.
\end{pitfall}

\subsection*{Repetition and columns}

A repetition can be computed while the program runs: in
\code|n (k) (...)|, the part in parentheses is repeated \code|k| times,
which suits arrays of any size. An integer value is printed through a
real picture as a real number:

\example{ch09/format-table}

The insertion \code|k| moves to the column given by its repetition
count, which can also be computed. That is enough for a simple plot:

\example{ch09/plot}

\noindent prints
\outputof{ch09/plot}

\subsection*{How the values and the format are matched}\label{sec-format-list}

The values of the list are matched with the pictures of the format one
by one. When the format is exhausted and values remain, it starts again
from the beginning: \code|printf (($2d l$, 1, 2, 3))| prints three lines.
This is the default action of the event \code|on_format_end|. A list may
contain several formats, each of which applies to the values after it.
The layout procedures \code|newline| and \code|space| cannot appear in
the list of \code|printf|; the insertions \code|l| and \code|x| take
their place.

A format is a value of the type \code|format| and can be stored in a
variable, passed to a procedure, or chosen at run time: a
\code|proc (int) format| can return different formats for different
cases. A format that computes a repetition from a local name,
\code|n (digits) d|, is subject to the lifetime rule
(Section~\ref{sec-proc-lifetimes}) like a function, and cannot leave the
block of \code|digits|.

\begin{implnote}
In the Report, the rest of a format that is not used up when the list
ends waits for the next \code|putf| on the same file. Algol 68 Genie
stops the program instead (``number of pictures does not match number
of arguments''): the values must use the format up to its end.
\end{implnote}

\section{Formatted input}\label{sec-readf}

\code|readf| and \code|getf| read according to a format. Every frame
reads one character position: a picture \code|3zd| reads four
characters, and a \code|z| frame accepts a blank as a zero. The
characters of a literal insertion such as \code|"-"| must be present in
the input and are read over; \code|x| skips one character. This makes it easy to read data in fixed columns, such as
the records of this file:
\lstinputlisting[style=output]{examples/ch09/readf.in}

\example{ch09/readf}

\noindent prints
\outputof{ch09/readf}

The pictures \code|"-" 2d| skip the dash and read the month. The
\code|10a| reads exactly ten characters, the blanks of a short name
included. A field that is one character too short or too long shifts
every field after it: with \code|4zd| for the score, which reads five
characters, the program stops at the first record with an error.

\section{Binary transput}\label{sec-binary}

The procedures \code|put_bin| and \code|get_bin| write values to a file
and read them back in an internal, binary form, without conversion into
characters. A file written this way can only be read back by a program
that reads the same types in the same order:

\example{ch09/binary}

The channel \code|stand_back_channel| allows both reading and writing,
and \code|reset| returns to the start of the file before reading. Binary
and character transput cannot be mixed on one sequential file until it
is reset; Algol 68 Genie enforces this.

Binary transput was meant for intermediate files that a program writes
and reads again. Such a file does not need a name, and it should
disappear when the program no longer needs it. \code|create| and
\code|scratch| do exactly that:

\example{ch09/temp}

\section{What Algol 68 Genie lacks}\label{sec-transput-genie}

Algol 68 Genie provides nearly all of the Report's transput. The
exceptions, apart from those already mentioned, are:
\begin{itemize}
\item the position enquiries \code|char_number|, \code|line_number| and
  \code|page_number|, \code|set_char_number|, \code|reidf|,
  \code|make_conv| and the event \code|on_char_error|;
\item \code|set| moves relative to the current position, with one
  parameter, instead of to a given page, line and character;
\item the insertion \code|k| cannot move back to an earlier column, and
  \code|backspace| and the insertion \code|y| cannot be used on
  \code|stand_out|;
\item a \code|z| frame after a \code|d| frame always prints its digit,
  where the Report suppresses some zeros in date-like pictures; pictures
  with only \code|d| frames after the first are printed the same
  everywhere;
\item the format is not kept for the next call of \code|putf|
  (Section~\ref{sec-format-list}).
\end{itemize}

\section{A complete program}\label{sec-transput-complete}

The final program reads a file of names and marks and prints a report.
The file \texttt{marks.txt} contains:
\lstinputlisting[style=output]{examples/ch09/marks.txt}

\example{ch09/report}

\code|make_term (data, " ")| ends a name at the first blank, so that
\code|get| reads the name and then the three numbers. The \code|newline|
at the end of the list skips the line break, so that the next name starts
at the beginning of the next line. The expression
\code|" " * (10 - UPB name)| pads the name to ten characters.

\section*{Exercises}
\addcontentsline{toc}{section}{Exercises}

The \hyperref[ans-ch9]{answers} are at the end of the book.

\begin{enumerate}
\item What is the difference between \code|open| and \code|establish|?
\item What does an event routine that returns \code|true| tell the
  transput?
\item Write a procedure \code|int_to_string| that converts an
  \code|int| into a string by writing it into a string with
  \code|associate| and \code|put|.
\item What do the formats \code|$4d$|, \code|$4zd$| and \code|$-3zd$|
  print for the values 7 and $-7$?
\item Why does the program of Section~\ref{sec-stand-in} read the line
  and the line break with two calls of \code|get|?
\item How many values does each of these formats print, and what do they
  print for the values 3 and 4?
  (a) \code|$2d "/" 2d$| (b) \code|$2d, "/" 2d$|
\item What do \code|whole (-5, 3)|, \code|whole (5, -3)|,
  \code|fixed (2.5, 0, 2)| and \code|fixed (-0.25, 0, 2)| yield?
\item A value error routine returns \code|true| without reading anything.
  What happens when \code|get (f, (a, b))| meets the input
  \code|x 5|?
\item Write a format that prints a date given as three integers day,
  month and year as \code|15 March 2024|.
\end{enumerate}

\section*{Programming problems}
\addcontentsline{toc}{section}{Programming problems}

\begin{enumerate}
\item \textbf{Line numbers.} Copy a text file to another file with a
  line number in front of each line, and start a new page with a page
  number after every 50 lines.
\item \textbf{Word count.} Count the characters, words and lines of the
  standard input, like the Unix program \texttt{wc}. Make sure that a
  last line without a line break is counted.
\item \textbf{Exam results.} Each line of a file holds a student
  number of 6 digits, a name in 20 columns and three marks of 3 digits.
  Print a line per student with the name, the average with one decimal,
  and \texttt{passed} or \texttt{failed} (an average of at least 50 and no
  mark below 30 passes) chosen with a \code|b(...)| pattern, with one
  \code|getf| and one \code|printf| per student.
\item \textbf{Column totals.} Read a table of numbers separated by
  commas, with an unknown number of rows, and print it with a format,
  with a total for each column.
\item \textbf{Bad records.} Read records of a fixed layout with
  \code|getf|, and install an \code|on_value_error| routine that prints
  the number of each bad record and skips it.
\item \textbf{Bank statement.} Read transactions (date, text, amount) and
  print a statement with a running balance, the amounts in a column with
  a floating sign or with \code|b("CR", "DR")|, and a heading on every
  page.
\item \textbf{A calendar.} Print the calendar of a month, with the name
  of the month chosen by a \code|c| pattern and the first week indented
  with a computed repetition.
\item \textbf{Amounts in words.} Print an amount of money below one
  million in words, as on a check, using choice patterns for the units,
  tens and hundreds.
\item \textbf{Justification.} Read a text and print it with lines of a
  given width and straight margins on both sides, by distributing extra
  blanks between the words as evenly as possible. The last line of a
  paragraph keeps a ragged right margin.
\item \textbf{The last lines.} Print the last $n$ lines of the standard
  input, which may be of any length, numbered with \code|whole|.
\item \textbf{Merging.} Write two sorted sequences of numbers to
  temporary files made with \code|create|, \code|reset| them, merge them
  onto the standard output, and \code|scratch| them.
\item \textbf{Number conversion.} Write a procedure that converts a
  \code|real| into a string with a given number of significant digits,
  using \code|associate| or \code|float|.
\item \textbf{Announcements.} Given the children of a family with
  their names, sexes and years of birth, print sentences such as ``Ann is the
  eldest of three children'' or ``Tom is the only son'', choosing the
  words with \code|b(...)|, \code|c(...)| and \code|f(...)| patterns in a
  single \code|printf| per child (Algol 68 Genie).
\end{enumerate}
